\documentclass[a4,11pt]{aleph-notas}
% Se puede ver la documentación aquí: 
% https://github.com/alephsub0/LaTeX_aleph-notas 

% -- Paquetes adicionales 
\usepackage{enumitem}
\usepackage{url}
\usepackage{array}
\usepackage{booktabs}
\usepackage{longtable}
\hypersetup{
    urlcolor=blue,
    linkcolor=blue,
}


% -- Datos
\institucion{Facultad de Ciencias Exactas, Naturales y Ambientale}
\carrera{Catálogo STEM}
\asignatura{Matemática III}
\tema{Plan tema 01: Matrices}
\autor{Andrés Merino}
\fecha{Periodo 2026-01}

\logouno[0.16\textwidth]{Logos/logoPUCE_04_ac}
\definecolor{colortext}{HTML}{1957d1}
\definecolor{colordef}{HTML}{1957d1}
\fuente{montserrat}


\begin{document}

\encabezado

%%%%%%%%%%%%%%%%%%%%%%%%%%%%%%%%%%%%%%%%  
\section*{Resultado de Aprendizaje}  
%%%%%%%%%%%%%%%%%%%%%%%%%%%%%%%%%%%%%%%%  

%%%%%%%%%%%%%%%%%%%%%%%%%%%%%%%%%%%%%%%%  
\subsection*{RdA de la asignatura:}  
% Comentar los que no apliquen a la clase
\begin{itemize}[leftmargin=*]  
    \item \textbf{RdA 1:} Identificar los conceptos, teoremas y propiedades de funciones vectoriales, derivadas parciales e integral vectorial.
    \item \textbf{RdA 2:} Calcular derivadas parciales, integrales de funciones vectoriales con el apoyo de asistentes computacionales.
    \item \textbf{RdA 3:} Aplicar distintos tópicos del Cálculo Vectorial para la solución de problemas reales en diferentes contextos.
\end{itemize}  

%%%%%%%%%%%%%%%%%%%%%%%%%%%%%%%%%%%%%%%%  
\subsection*{RdA del tema:}  
% indicar de 2 a 3 resutados de aprendizaje específicos para este tema
\begin{itemize}[leftmargin=*]  
    \item Identificar matrices según su orden, filas y columnas.  
    \item Aplicar operaciones básicas entre matrices: suma, producto por escalar y transpuesta.  
    \item Determinar cuándo una multiplicación de matrices está definida y calcular el producto.  
\end{itemize}  

%%%%%%%%%%%%%%%%%%%%%%%%%%%%%%%%%%%%%%%%  
\section*{Introducción}  
%%%%%%%%%%%%%%%%%%%%%%%%%%%%%%%%%%%%%%%%  

%%%%%%%%%%%%%%%%%%%%%%%%%%%%%%%%%%%%%%%%  
\paragraph{Pregunta inicial:}  
% Colocar una pregunta que motive desde la cotidianidad o la ciencia de datos el tema a tratar, que no involucre en el enunciado los concpetos que se van a abordar.
¿Cómo se puede organizar información de manera compacta para modelar problemas de ingeniería?

%%%%%%%%%%%%%%%%%%%%%%%%%%%%%%%%%%%%%%%%  
\section*{Desarrollo}  
%%%%%%%%%%%%%%%%%%%%%%%%%%%%%%%%%%%%%%%%  

%%%%%%%%%%%%%%%%%%%%%%%%%%%%%%%%%%%%%%%%  
\subsection*{Actividad 1: Modelación con tablas y representación matricial}
% Breve descripción
La clase inicia con una tabla de datos reales (por ejemplo: producción semanal por sedes y productos) para motivar la noción de matriz como arreglo numérico. A través de una clase magistral, se trabajan los conceptos de orden, igualdad de matrices, operaciones elementales y producto matricial, enfatizando la compatibilidad de dimensiones y su interpretación.

\paragraph{¿Cómo lo haremos?}  
% indicar las que correspondan
\begin{itemize}[leftmargin=*]  
    \item \textbf{Motivación:} se presenta una tabla de datos y se discute cómo convertirla a matriz, identificando filas, columnas y entradas $a_{ij}$.
    \item \textbf{Clase magistral:} se explican los conceptos de matriz, orden, igualdad, suma, multiplicación por escalar, transpuesta, producto de matrices y matriz identidad. Se usa el resumen \href{https://fcena-puce.github.io/AlgebraLineal-05-N0048/01Resumen/Resumen01.pdf}{Resumen01.pdf}.  
    \item \textbf{Resolución de ejercicios:} los estudiantes resolverán, de manera guiada, ejercicios de operaciones y producto de matrices del documento \href{https://fcena-puce.github.io/AlgebraLineal-05-N0048/04Ejercicios/Ejercicios01.pdf}{Ejercicios01.pdf}.  
    \item \textbf{Visualización de videos:} se recomendará el estudio con los videos:  
    \begin{itemize}
        \item \href{https://youtu.be/9FKFgNQktkU?si=kGRbP9Ad-NRRciwL}{MATRICES: de los gráficos de Fortnite a la física cuántica}
        \item \href{https://youtu.be/m6w5vLA3Lnw?si=xcCeg-0snioQb-ns}{Matrices Introducción | Conceptos básicos}
    \end{itemize}  
    \item \textbf{Recomendaciones de lectura:} se sugiere la lectura de los siguientes enlaces:
    \begin{itemize}
        \item \href{https://blog.nekomath.com/algebra-lineal-i-bases-y-dimension-de-espacios-vectoriales/}{Bases y dimensión de espacios vectoriales}
    \end{itemize}
    \item \textbf{Implementación computacional:} se realiza mediante cuadernos de Jupyter usando el documento \href{https://fcena-puce.github.io/AlgebraLineal-05-N0048/02ResumenPython/ResumenPy01.pdf}{ResumenPy01.pdf}.
\end{itemize}  

\paragraph{Verificación de aprendizaje:} 
% plantear preguntas directas y cerradas que permitan verificar que se han alcanzado los RdA del tema
\begin{itemize}[leftmargin=*]  
    \item ¿Qué condiciones deben cumplir dos matrices para poder sumarse?  
    \item ¿Cuándo está definido el producto $AB$ entre dos matrices?  
    \item ¿Por qué, en general, $AB \neq BA$?  
\end{itemize}  

%%%%%%%%%%%%%%%%%%%%%%%%%%%%%%%%%%%%%%%%  
\section*{Cierre}  
%%%%%%%%%%%%%%%%%%%%%%%%%%%%%%%%%%%%%%%%  

\paragraph{Tarea:}  
Resolver del libro \href{https://catalogobiblioteca.puce.edu.ec/cgi-bin/koha/opac-detail.pl?biblionumber=86081}{Fundamentos de álgebra lineal de Larson}, sección 2.1, los ejercicios: 2, 4, 10, 11, 14, 18, 20, 31, 35, 62, 63, 81.

\paragraph{Pregunta de investigación:}  
% preguntas relacionadas con los temas que no se han abordado en clase
\begin{enumerate}[leftmargin=*]  
    \item ¿Cómo se usan las matrices para representar imágenes digitales y transformarlas (rotación, escalamiento, filtros)?
    \item ¿Qué relación existe entre matrices y redes neuronales en inteligencia artificial?
\end{enumerate}  

\paragraph{Para el próximo tema:}  
Leer la sección 2.3 del libro \href{https://catalogobiblioteca.puce.edu.ec/cgi-bin/koha/opac-detail.pl?biblionumber=283260}{Álgebra lineal con aplicaciones y Python de Aranda} y la sección 1.1 del libro \href{https://catalogobiblioteca.puce.edu.ec/cgi-bin/koha/opac-detail.pl?biblionumber=86081}{Fundamentos de álgebra lineal de Larson}.

\end{document}
